\documentclass[10pt,a4paper]{article}
\usepackage[margin=15mm,top=12mm,bottom=12mm]{geometry}
\usepackage{fontspec}
\usepackage[french]{babel}
\usepackage{amsmath,amssymb}
\usepackage{xcolor}
\usepackage[most]{tcolorbox}
\usepackage{booktabs}
\usepackage{enumitem}
\usepackage{graphicx}
\usepackage{array}
\usepackage{titlesec}
\usepackage{listings}
\usepackage{hyperref}

\definecolor{bleu}{HTML}{1F6FEB}
\definecolor{gris}{HTML}{5B6470}
\definecolor{vert}{HTML}{1E8449}
\definecolor{orange}{HTML}{B45309}
\definecolor{fond}{HTML}{F4F6F8}
\definecolor{codebg}{HTML}{F6F8FA}

\setlength{\parindent}{0pt}
\titleformat{\section}{\normalfont\bfseries\color{bleu}\large}{}{0pt}{}
\titlespacing*{\section}{0pt}{9pt}{4pt}
\titleformat{\subsection}{\normalfont\bfseries\normalsize\color{bleu}}{}{0pt}{}
\titlespacing*{\subsection}{0pt}{7pt}{2pt}
\setlist[itemize]{nosep,leftmargin=1.1em,labelsep=.45em}
\setlist[enumerate]{nosep,leftmargin=1.35em,labelsep=.5em}

\lstdefinestyle{py}{language=Python,backgroundcolor=\color{codebg},
  basicstyle=\ttfamily\fontsize{7.6}{9}\selectfont,
  keywordstyle=\color{bleu}\bfseries,commentstyle=\color{vert}\itshape,
  stringstyle=\color{orange},frame=none,breaklines=true,
  showstringspaces=false,columns=fullflexible,keepspaces=true,
  aboveskip=2pt,belowskip=2pt,
  literate={é}{{\'e}}1 {è}{{\`e}}1 {à}{{\`a}}1 {ç}{{\c c}}1 {û}{{\^u}}1}
\lstset{style=py}

\newtcolorbox{obj}{enhanced,colback=fond,colframe=fond,arc=2pt,
  left=6pt,right=6pt,top=4pt,bottom=4pt,breakable}
\newtcolorbox{note}{enhanced,colback=orange!5,colframe=orange!65!black,boxrule=.5pt,
  arc=2pt,left=6pt,right=6pt,top=4pt,bottom=4pt,breakable}
\newtcolorbox{bareme}{enhanced,colback=bleu!4,colframe=bleu!55!black,boxrule=.5pt,
  arc=2pt,left=6pt,right=6pt,top=4pt,bottom=4pt,breakable}

\begin{document}

% ===================================================== EN-TETE
\begin{tcolorbox}[enhanced,colback=fond,colframe=fond,arc=3pt,
  left=9pt,right=9pt,top=7pt,bottom=7pt]
\begin{minipage}[c]{0.76\textwidth}
  {\color{gris}\footnotesize LICENCE 3 SDIA $\cdot$ ENSPD --- UNIVERSITÉ DE DOUALA}\\[2pt]
  {\Large\bfseries Travaux pratiques --- énoncés TP1 et TP2}\\[3pt]
  {\small Introduction au Machine Learning $\cdot$ 60 h $\cdot$ 2 séances de 2 h}\\[4pt]
  {\color{gris}\footnotesize Dr.\ Noulapeu Ngaffo Armielle}
\end{minipage}\hfill
\begin{minipage}[c]{0.20\textwidth}
  \raggedleft\includegraphics[width=17mm]{assets/logo_enspd.png}
\end{minipage}
\end{tcolorbox}

\begin{note}
\textbf{Modalités communes.}
Chaque TP se fait \textbf{seul}, sur machine, en \textbf{2 h}.
Le rendu est le \textbf{notebook exécuté} (toutes les cellules, sorties visibles)
déposé sur la plateforme du cours avant la séance suivante.
La graine est fixée à \texttt{4112} : deux exécutions doivent donner exactement
les mêmes nombres. Tout tourne sur CPU, sans téléchargement.
Les cellules marquées \textbf{« À vous »} sont les seules notées ; les autres
sont guidées et servent de point de départ.
\end{note}

% ===================================================== TP1
\section{TP1 --- Premier pipeline et première fuite}
{\color{gris}\small Durée : 2 h $\cdot$ Barème : 10 points $\cdot$
Jeu : \texttt{breast\_cancer} (569 patients, 30 mesures, bénin / malin)}

\begin{obj}
\textbf{Objectifs.}
\begin{enumerate}
  \item Charger des données réelles et les découper correctement.
  \item Comparer un modèle à une \textbf{baseline} : sans elle, aucun score n'est interprétable.
  \item Lire une matrice de confusion et une validation croisée.
  \item Fabriquer une fuite de données, puis mesurer son effet.
\end{enumerate}
\end{obj}

\subsection*{Déroulé guidé (non noté)}
\begin{enumerate}
  \item \textbf{Charger} le jeu et vérifier les dimensions et l'équilibre des classes
        (\texttt{value\_counts}).
  \item \textbf{Découper} en entraînement et test (\texttt{test\_size=0{,}3},
        \texttt{random\_state=4112}, \texttt{stratify=y}).
  \item \textbf{Baseline} : prédire toujours la classe majoritaire ; mesurer
        l'exactitude. C'est le chiffre à battre.
  \item \textbf{Pipeline} : standardisation puis régression logistique, ajustés
        ensemble (\texttt{Pipeline}) ; mesurer sur le test.
  \item \textbf{Matrice de confusion} : identifier les faux négatifs et les faux
        positifs, puis calculer précision et rappel.
  \item \textbf{Validation croisée} à 5 plis : rapporter la moyenne \emph{et}
        l'écart-type.
  \item \textbf{Fabriquer une fuite} : ajouter 300 colonnes de bruit pur et tirer
        les étiquettes au hasard. Sélectionner les variables \emph{avant} la
        validation croisée et observer le score. Refaites la sélection
        \emph{dans} chaque pli et comparez.
\end{enumerate}

\subsection*{Travail noté}
\begin{bareme}
\textbf{Question 1 (4 points).}
Remplacez la régression logistique par un \texttt{RandomForestClassifier}
(\texttt{n\_estimators=200}, \texttt{random\_state=4112}).
La baseline change-t-elle ? Le gain par rapport à la baseline est-il plus grand
ou plus petit qu'avec la régression logistique ? Concluez en une phrase.
\end{bareme}

\begin{bareme}
\textbf{Question 2 (3 points).}
Le jeu est-il équilibré ? Regardez \texttt{y.value\_counts()}.
Quel serait le danger de juger ce modèle avec l'exactitude seule ? Donnez le
chiffre de la classe minoritaire et expliquez en deux phrases.
\end{bareme}

\begin{bareme}
\textbf{Question 3 (3 points).}
Dans l'étape 7, les étiquettes sont \emph{aléatoires}. Expliquez pourtant, en
deux phrases, pourquoi la sélection faite avant la validation croisée produit un
score élevé, et pourquoi la sélection faite dans chaque pli ramène le score au
niveau du hasard.
\end{bareme}

\subsection*{Ce qui est attendu dans le rendu}
\begin{itemize}
  \item Les trois réponses rédigées dans des cellules markdown, avec les
        \textbf{nombres obtenus} cités.
  \item Le notebook exécuté de bout en bout, sans erreur.
  \item Une phrase de conclusion personnelle sur ce que la baseline a changé
        dans votre lecture du score.
\end{itemize}

\newpage
% ===================================================== TP2
\section{TP2 --- Préparer les données sans fuite}
{\color{gris}\small Durée : 2 h $\cdot$ Barème : 10 points $\cdot$
Jeu : synthétique (400 dossiers, graine 4112)}

\begin{obj}
\textbf{Objectifs.}
\begin{enumerate}
  \item Diagnostiquer les valeurs manquantes, les types et les échelles.
  \item Construire un \texttt{ColumnTransformer} complet.
  \item Ajouter un encodage cyclique pour une variable horaire.
  \item Fabriquer une fuite par encodage cible, puis la corriger.
\end{enumerate}
\end{obj}

\subsection*{Déroulé guidé (non noté)}
\begin{enumerate}
  \item \textbf{Diagnostiquer} : nombre de manquants par colonne, types, valeurs
        extrêmes. On ne corrige que ce que l'on a vu.
  \item \textbf{Découper avant toute transformation} : le test ne doit jamais
        servir à calculer une moyenne, une médiane ou un encodage.
  \item \textbf{Pipeline naïf} : imputation par la médiane et encodage one-hot,
        ajustés sur l'entraînement seulement ; mesurer.
  \item \textbf{Encodage cyclique} de l'heure : deux colonnes $\sin$ et $\cos$
        pour que 23 h et 0 h soient voisines ; comparer.
  \item \textbf{Fabriquer une fuite} : variable « agence » à 100 modalités,
        2 lignes chacune, étiquette tirée au hasard. Calculer l'encodage cible
        sur \emph{tout} le jeu et observer le score.
  \item \textbf{Corriger} : recalculer l'encodage à l'intérieur de chaque pli
        (ou utiliser \texttt{TargetEncoder} dans le pipeline) et comparer.
\end{enumerate}

\subsection*{Travail noté}
\begin{bareme}
\textbf{Question 1 (4 points).}
Remplacez l'imputation médiane du revenu par une imputation par la
\textbf{moyenne}. Le score change-t-il ? De combien de points ? Expliquez
pourquoi l'écart est faible ou nul, en vous appuyant sur ce que vous avez
observé dans le diagnostic (asymétrie, valeurs extrêmes).
\end{bareme}

\begin{bareme}
\textbf{Question 2 (3 points).}
Retirez la valeur aberrante (\texttt{9\,000\,000}) avant l'entraînement.
Le rappel s'améliore-t-il ? Donnez les deux valeurs (avant / après) et concluez
sur l'effet d'une seule observation sur un jeu de 400 lignes.
\end{bareme}

\begin{bareme}
\textbf{Question 3 (3 points).}
Dans l'étape 5, l'étiquette est tirée au hasard. Expliquez en deux phrases
pourquoi l'encodage cible calculé sur tout le jeu produit un score élevé, et
pourquoi le recalcul par pli le ramène au hasard.
\end{bareme}

\subsection*{Ce qui est attendu dans le rendu}
\begin{itemize}
  \item Les trois réponses rédigées, avec les \textbf{nombres obtenus}.
  \item Le notebook exécuté de bout en bout, sans erreur.
  \item Le tableau comparatif « pipeline naïf / avec encodage cyclique /
        encodage cible fuité / encodage cible corrigé ».
\end{itemize}

\vspace{2mm}
\begin{note}
\textbf{Barème commun aux deux TP.}
Exactitude des calculs et des chiffres cités : 6 points.
Justification en deux phrases, vocabulaire du cours : 3 points.
Notebook exécuté sans erreur et lisible : 1 point.\\[2pt]
Un résultat faux mais une démarche correcte et explicitée vaut mieux qu'un
résultat juste non justifié.
\end{note}

\end{document}
